% ============================================================
% Section 5: Discussion
% ============================================================
\section{Discussion}
\label{sec:discussion}

\subsection{Why a Probabilistic Phase Transition?}
\label{sec:why_probabilistic}

The S-shaped $P(\text{success} \mid N, k)$ curves documented in
\S\ref{sec:success_curves} merit an explanation.  Why should the same
experiment---same model, same data, same optimizer---produce either random
guessing or near-perfect classification depending only on the random seed?
And why does this behavior persist over a wide range of $N$ rather than
collapsing to a sharp threshold?

\paragraph{A geometric intuition.}
Consider the empirical loss landscape $L_N(\mathbf{w})$ induced by $N$
CRPs.  The true weight vector $\mathbf{w}^*$ defines a global minimum
corresponding to the exact PUF response function.  When $N$ is small,
$L_N$ is a poor approximation of the population loss: the terrain is flat,
with many spurious minima that all yield $\approx 50\%$ accuracy.  As $N$
grows, a \emph{basin of attraction} forms around $\mathbf{w}^*$, and the
optimizer (SGD) may fall into it if randomly initialized within its
catchment area.

The probability $P(\text{success})$ is therefore the probability that a
random initialization lies in the basin of $\mathbf{w}^*$.  This basin grows
with $N$:
\[
\text{Vol}(\text{basin}(\mathbf{w}^*; N)) \propto
\Phi\!\left(\frac{N - N_c}{\sigma_N}\right),
\]
where $\Phi$ is a sigmoidal function, $N_c$ is the scaling midpoint, and
$\sigma_N$ governs the transition width.  The S-curve is thus a geometric
consequence of a growing attractor volume in a high-dimensional parameter
space.

\paragraph{A heuristic derivation of the exponential scaling.}
A soft-bit relaxation of the XOR parity makes the information bottleneck
precise.  Write the delay of the $j$-th APUF as
$D_j = \mathbf{w}_j \cdot \mathbf{c} + b_j$ and the soft bit
$s_j = \sigma(D_j)$.  The soft response probability is
$P(r{=}1) = \frac{1}{2}\bigl(1 - \prod_{j=1}^{k}(1-2s_j)\bigr)$, which
recovers the exact parity in the hard limit.  The log-likelihood gradient
with respect to $D_i$ is
\[
\frac{\partial \log P}{\partial D_i} \;=\;
s_i(1-s_i) \cdot \prod_{j \neq i}(1-2s_j),
\]
up to a bounded factor depending on the label $r$.  The first factor
$s_i(1-s_i)$ is $\Theta(1)$ for generic delays, but the product over the
\emph{other} $k-1$ soft bits is the attenuation: for random parameters,
$s_j$ is symmetrically distributed about $1/2$, giving
$\mathbb{E}|1-2s_j| \leq 1/2$, and hence
\[
\mathbb{E}\Bigl|\prod_{j \neq i}(1-2s_j)\Bigr|
\;\approx\; \bigl(\mathbb{E}|1-2s_j|\bigr)^{k-1}
\;\lesssim\; 2^{-(k-1)}.
\]
Each CRP therefore carries an expected gradient signal of order
$2^{-(k-1)}$ toward every delay parameter.  Accumulating an $\Omega(1)$
total signal---the minimum required for the optimizer to escape the flat
region around the trivial all-zero solution---requires
$N_c \gtrsim 2^{k-1}$ CRPs, i.e., \emph{exponential growth with base 2 per
additional XOR level}.  Our measurements give per-level factors of
$3.4$--$6.1$, above this lower bound: the excess reflects optimizer
inefficiency and the ruggedness of the empirical landscape, both of which
compound with $k$.  The same attenuation governs the likelihood surface of
the 520-dimensional delay parameterization: candidate pairs differing in a
single delay produce NLL gaps of order $2^{-(k-1)}$, so evolution strategies
(\S\ref{sec:evolution_baselines}) inherit the identical data threshold even
though their global optimum is the true parameter vector.  This is a
heuristic argument, not a theorem; the constants governing
$N_c$ remain uncontrolled, and the precise geometry of the basin boundary
is an open question.

\paragraph{Why not a sharp threshold?}
If the basin volume were a step function, $P(\text{success})$ would jump
from 0 to 1 at a critical $N_c$.  In practice, the basin boundary is fractal
in high dimensions~\cite{choromanska2015loss}, and random initializations
sample its volume continuously.  The transition width varies
nonmonotonically with $k$ in our data.  Higher-XOR landscapes are more
rugged~\cite{dauphin2014identifying}: the basin of $\mathbf{w}^*$ is smaller
and more isolated, so a wider absolute range of $N$ is needed for random
initializations to intersect it.  This explains the widening S-curve in
absolute terms from 2-XOR (transition spans $10^4$--$10^5$ CRPs) to 8-XOR
(transition spans $5\times10^6$--$10^7$ CRPs) to 9-XOR (transition likely
exceeds $1.5\times10^7$ CRPs).

\paragraph{Relation to grokking and phase transitions in deep learning.}
Our observed phenomenon shares qualitative features with ``grokking'' in
overparameterized networks~\cite{power2022grokking} and phase transitions in
neural network learning~\cite{saxe2014dynamics}.  In grokking, validation
accuracy remains near random for many epochs before suddenly rising to
perfect generalization; here, the transition occurs along the data axis
rather than the epoch axis.  Both phenomena reflect a delayed emergence of
structured solutions in high-dimensional nonconvex optimization.  The
probabilistic element---that some runs succeed while others fail at the same
$N$---is less commonly reported and may be specific to the XOR APUF problem,
where the target function is known to be exactly realizable by the model.

\subsection{Effect of Physical Noise on the Phase Transition}

All experiments in this work use idealized, noise-free PUF simulations.
Real silicon PUFs introduce thermal noise, supply-voltage fluctuations, and
measurement jitter that cause a fraction of CRP bits to be unreliable.
We briefly outline the expected effect on the S-curve framework.

Consider a PUF whose noise-free response $y \in \{0,1\}$ is corrupted by
independent bit-flip noise with probability $\varepsilon$.  The mutual
information between the noisy and clean responses is
$I = 1 - H(\varepsilon)$, where $H(\varepsilon) = -\varepsilon\log_2\varepsilon
- (1-\varepsilon)\log_2(1-\varepsilon)$ is the binary entropy.  The effective
information per CRP is therefore reduced by a factor of approximately
$1 - H(\varepsilon)$.  This suggests a simple rescaling: the CRP budget
required to achieve a given attack success probability would increase by
roughly $1/(1-H(\varepsilon))$.

A preliminary experimental validation supports this prediction.  We trained
the standard MLP on 6-XOR with 400{,}000 CRPs under bit-flip noise levels
$\varepsilon \in \{0, 3\%, 5\%\}$, with evaluation on clean test data.  The
observed success probabilities were 80\% (4/5), 60\% (3/5), and 40\% (2/5),
respectively (results/noise\_6xor\_400k\_v2.csv).  The noiseless baseline
exactly matches the earlier 8/10 = 80\% success probability at this $(k, N)$ point
(baseline\_mlp\_lr.csv), confirming protocol consistency.  The monotonic
decrease in success probability with noise level is consistent with the
information-theoretic rescaling prediction, though the sample size (5 runs
per condition) permits only a qualitative comparison.

A complementary measurement at higher XOR complexity shows a stronger
effect.  At 8-XOR with 5{,}000{,}000~CRPs, $\varepsilon = 3\%$ noise reduced
the empirical success probability from 28\% (7/25, noiseless) to 0\% (0/10;
results/noise\_8xor\_5M.csv)---consistent with the predicted rightward shift
of the S-curve at higher XOR levels.  The effect appears stronger at $k=8$
than at $k=6$ (80\% $\to$ 60\% at 400k), suggesting that the noise-induced
shift compounds with XOR complexity, though the sample sizes (5--10 per
condition) preclude quantitative comparison.

Under typical environmental noise levels of $\varepsilon = 3$--$5\%$, this
factor is between approximately $1.24$ and $1.40$ (a hypothetical estimate
based on an idealized noise model, not an experimental measurement)---a
rightward shift of the S-curve by approximately $24$--$40\%$ in CRP budget.  The qualitative phenomena---the existence of a
probabilistic transition and the architecture independence---are preserved,
as the key determinant remains the information-theoretic data requirement
rather than the model family.

A more subtle effect is that nonuniform noise (e.g., bias-temperature
instability causing systematic errors on specific APUF stages) could
distort the S-curve shape rather than purely translating it.  Modeling
these effects---via simulation with injected noise or through silicon
measurements---is a natural extension of this work.  A preliminary
simulation at 6-XOR/400k with $\varepsilon \in \{3\%, 5\%\}$ confirms the
predicted direction (success probability 80\% $\to$ 60\% $\to$ 40\% as
$\varepsilon$ increases from 0 to 5\%); a comprehensive noise grid across
$(k, N)$ configurations and a silicon validation remain future work.

\subsection{Implications for PUF Security Evaluation}
\label{sec:implications}

Our findings have direct consequences for how the security of strong PUFs
should be evaluated and reported.

\paragraph{Security is probabilistic, not binary.}
The standard practice of declaring a PUF design ``broken'' when a single
attack demonstration succeeds is misleading.  For the 8-XOR APUF, one can
demonstrate either 99\% accuracy (with 10M CRPs and the right seed) or 50\%
accuracy (with the same 10M CRPs and a different seed).  A responsible
evaluation must report $P(\text{success} \mid N, k)$ across multiple trials,
along with the CRP budget.  We recommend that future work adopt the S-curve
framework as a standard reporting format.

\paragraph{$k$-XOR security quantification.}
Based on our data, the CRP budget required for an attack scales by
roughly 3--6$\times$ per additional XOR level for $k=6\to7\to8$ (the only
transitions with reliable $N_{50}$ estimates); the $k=2\to4$ transition is
anomalous ($N_{50}$ decreases):
\begin{itemize}
    \item 2-XOR: $\sim$103{,}000 CRPs ($N_{50}$, logistic fit);
    \item 4-XOR: $<$75{,}000 CRPs ($N_{50}$; no data below 75k, estimate unreliable);
    \item 6-XOR: $\sim$384{,}000 CRPs ($N_{50}$, logistic fit);
    \item 7-XOR: $\sim$1{,}320{,}000 CRPs ($N_{50}$, logistic fit);
    \item 8-XOR: $\sim$8{,}000{,}000 CRPs ($N_{50}$, logistic fit, 95\% CI [6.2M,~9.8M]);
    \item 9-XOR: $>1.5\times10^7$ CRPs (very rough lower bound; 0/3 successes at 15M with CI [0\%,71\%]; 29 total trials across 5 CRP levels).
\end{itemize}
An 8-XOR APUF with CRP exposure limited to $5 \times 10^6$ provides
practical security against single-GPU ML attacks on our test hardware
(NVIDIA RTX 3080 Ti, 12~GB VRAM): the attacker faces at best
a 28\% success probability even with optimal model choices.  For 9-XOR,
the preliminary $N_{50}$ estimate of $>10^7$ CRPs already exceeds the memory
capacity of our 12~GB GPU; measurements at 12{,}000{,}000~CRPs were possible
only on a 24~GB RTX~4090.  At $k = 10$, $N_{50}(10)$ exceeds $10^7$ CRPs
(0/3 successes at 10{,}000{,}000~CRPs), with the transition point
unobservable on current hardware: attempts at 15{,}000{,}000~CRPs fail with
out-of-memory errors even on 24~GB VRAM.  Any quantitative bound beyond
$10^7$ is therefore an engineering inference, not an empirical estimate.
Below the transition, even evolution strategies with the correct model
structure fail; the data threshold is a property of the information
content, not of the optimizer.

\paragraph{Design recommendations.}
Our measurements translate into concrete guidance for PUF-based key
generation and authentication:
\begin{itemize}
    \item \textbf{8-XOR with a lifetime CRP cap.}  An 8-XOR APUF whose
    challenge interface releases fewer than $5\times10^6$ CRPs over its
    lifetime retains a comfortable margin against single-GPU ML attacks:
    at the cap the attacker succeeds in at most 28\% of trials, and we
    observed no successful run at or below $2\times10^6$ CRPs.  Moving the
    cap to $5.5\times10^6$ confers no benefit---the success probability rises to 50\%.
    \item \textbf{9-XOR for higher assurance margins.}  A 9-XOR APUF shows
    no transition up to $1.5\times10^7$ CRPs, beyond the largest attack we
    could mount on a 24~GB GPU.  The marginal hardware cost of one
    additional XOR stage is small relative to the $>5\times$ increase in
    data complexity.
    \item \textbf{Silicon noise is a free hardening layer.}
    $\varepsilon = 3\%$ bit-flip noise erases the residual 28\% success at
    8-XOR/$5\times10^6$ CRPs (0/10 in our grid), shifting the S-curve
    rightward at no design cost---provided the PUF's own error-correction
    requirements remain satisfiable.
    \item \textbf{Avoid iPUF designs.}  The split attack reduces the
    effective security of every tested iPUF configuration by exactly 2 XOR
    levels, while its interpose stage adds hardware cost.  A designer who
    needs $(k_{\text{up}}+k_{\text{down}})$-XOR security obtains it more
    cheaply---and without the structural leakage---from a plain XOR APUF of
    that size.
\end{itemize}

\paragraph{Choice of success criterion.}
We adopt 90\% test accuracy as the success criterion throughout this study.
Cryptographic key recovery can sometimes exploit partial information at lower
accuracy levels (60--85\%), so the 90\% threshold represents a conservative
bound on practical key extraction.  The qualitative S-curve shape is
insensitive to this choice: shifting the threshold to 80\% or 95\% shifts
$N_{50}$ by less than 10\% for all $k$, preserving the relative ordering
across XOR levels.

\paragraph{Architecture does not matter; data does.}
The architecture independence result (\S\ref{sec:arch_independence}) has a
practical corollary: an attacker gains nothing from sophisticated model
design or hyperparameter tuning below the data threshold.  Conversely, a
defender cannot assume that architectural obscurity will protect against
attack---if sufficient CRPs are available, a standard MLP suffices.

\paragraph{iPUF is not a security solution.}
Our iPUF results (\S\ref{sec:ipuf}) show that the split architecture
dramatically reduces data complexity, and---beyond the known breakability
result---that the reduction is \emph{quantifiable}: the equivalent XOR
complexity drops by exactly 2 for every tested configuration, so the
CRP-budget gap against the claimed level grows from $\approx$5$\times$ at
$(2,4)$ to $\approx$21$\times$ at $(4,4)$ on our fitted scaling law.
Designers should therefore treat the claimed
$(k_{\text{up}}+k_{\text{down}})$ equivalence as an upper bound that is
systematically violated by 2 XOR levels in the attacker's favor.

\subsection{Limitations}
\label{sec:limitations}

We restrict our comparison to MLP-based architectures and two
evolution-strategy baselines (PSO, CMA-ES; \S\ref{sec:evo_baselines})
because prior work has
established that CMA-ES-based methods (Becker~2015,
CalyPSO~2024~\cite{calypso2024}) require comparable or larger CRP budgets
than MLPs for XOR APUFs of $k \geq 6$, and our focus is on demonstrating
that neither model architecture nor optimization paradigm has an effect
below the transition.

Several caveats apply to our findings.

First, our $N_{50}$ estimates for 2-XOR ($1.03\times10^5$) and 4-XOR
(below the lowest tested 75k CRP budget) are \emph{extrapolated}: the
2-XOR fit extrapolates above the highest tested budget (100k), while the
4-XOR fit extrapolates below the lowest tested budget (75k).  Additional
experiments at 60k--90k for 2-XOR and 50k--60k for 4-XOR would sharpen
these estimates considerably.

Second, the 6-XOR 300k data point shows high variability (33\% success,
24~trials) that may confound precise $N_{50}$ estimation; denser sampling
in the 200k--300k range would further resolve the shape of the transition
at this critical XOR level.

Third, our 9-XOR results are based on 29 total trials across five CRP
levels (5M, 7.5M, 10M, 12M, and 15M) and should be treated as preliminary.
The estimated $N_{50}(9) > 1.5\times10^7$ is subject to large uncertainty,
especially at 12M and 15M where only three trials were run at each level
(0/3 at both).

Fourth, the architecture independence claim has been verified only at
$(k=8, N=2\text{M})$ and one additional sub-threshold configuration;
testing at additional $(k,N)$ pairs---particularly for 6-XOR and 7-XOR
below their respective transitions---would strengthen the generalization.

Fifth, all experiments use simulated PUF models with Gaussian weight
distributions and zero measurement noise.  Silicon PUFs introduce thermal
noise, bias, and environmental variation, which may shift or widen the
observed transition.

Sixth, this study is limited to MLP architectures and two
evolution-strategy baselines under uniform random
challenge protocols.  Generalizability to other model families (CNNs,
transformers, kernel methods) and attack strategies (chosen-challenge,
adaptive-challenge~\cite{sayadi2025}) remains an open question.

Seventh, we tested only the original iPUF design (Nguyen~et~al.~2019);
obfuscation-enhanced or noise-hardened iPUF variants may exhibit different
vulnerability profiles.

Eighth, our computational budget is limited to two GPUs---an NVIDIA
RTX~3080 Ti (12~GB VRAM) and an RTX~4090 (24~GB VRAM)---which restricts
$k \geq 10$ to at most $10^7$ CRPs (0/3 successes) and prevents
15{,}000{,}000~CRPs entirely (out-of-memory).  Consequently, the
$N_{50}(10)$ transition point is unobservable on current hardware, and any
quantitative bound beyond $10^7$ CRPs is an engineering inference, not an
empirical estimate.  Future work with larger GPU memory could extend the
scaling law to higher XOR complexities.

These limitations suggest promising directions for future work but do not
diminish the central finding: for MLP and evolution-strategy attacks on XOR
APUFs under random challenges, attack success is a probabilistic function of
CRP budget, not a deterministic one.
